\uuid{Sz3c}
\titre{Noyau/Image/Injection/Surjection/Théorème du rang}
\auteur{hanani ; mzari}
\organisation{exo7}
\licence{CC BY-NC-SA 4.0}
\datecreate{2020-01-20}
\niveau{L1}
\module{Algèbre}
\chapitre{Application linéaire}
\sousChapitre{Définition}
\difficulte{3}
\theme{Noyau et image, Noyau, Image, Injection, Surjection, Théorème du rang}
\exo7qcmid{lille-2.applications-lineaires.023}

\contenu{
\begin{question}
\qclassification{107.01}
\qtags{motcle=Noyau/Image/Injection/Surjection/Théorème du rang}
Soit $E$ et $F$ deux $\Rr$-espaces vectoriels de dimensions finies et $f$ une application linéaire de $E$ dans $F$. On pose $\dim E=n $ et $\dim F=m$. Quelles sont les assertions vraies ?
\begin{answers}
\good{Si $f$ est injective, alors $n \le m$}.
\bad{Si $n \le m$, alors $f$ est injective}.
\good{Si $f$ est surjective, alors $n \ge m$}.
\bad{Si $n \ge m$, alors $f$ est surjective}.
\end{answers}
\begin{explanations} On a : $\dim Im f \le \dim F=m$.
\vskip0mm
Du théorème du rang, on déduit que si $n>m$,  $\dim \ker f >0$ et donc $f$ n'est pas injective et que si $n<m$,  $\dim \Im f <m$ et donc $f$ n'est pas surjective.
\vskip0mm
Si $n\le m$, alors $f$ n'est pas nécessairement injective et si $n \ge m$, alors $f$ n'est pas nécessairement surjective. Exemple : L'application nulle de  $\Rr$ dans $\Rr$.
\end{explanations}
\end{question}
}
